\section{从“求和”要求中辨认机器必须保留的事实}

先不画整机框图。设想处理器已经开始求和：它刚刚读完第一个整数 \(7\)，下一步准备读取
\(-2\)。如果此刻把机器中的全部状态抹去，再告诉它“继续”，它无法判断下一条指令在哪里，
也不知道已经累加出的部分和、尚未读取的数组位置以及剩余元素个数。因此，“运行”不能只
表示某段指令存在；运行意味着若干事实在相邻动作之间得以保留。

对这个求和循环，至少有五项事实不能丢失：

\begin{enumerate}
\item 下一条机器指令的地址；
\item 当前部分和；
\item 下一个输入元素的地址；
\item 尚未处理的元素个数；
\item 函数返回时应当回到的位置。
\end{enumerate}

这些事实的含义来自程序约定，电路只需要保存对应的位串。教学处理器把“下一条指令地址”
单独保存在程序计数器（program counter，PC）中，把计算中的位串保存在通用寄存器中，把
函数调用形成的返回地址保存在由栈指针（stack pointer，SP）定位的内存区域中。PC 和 SP
的名称描述用途，不表示里面的位具有特殊材料；它们仍是普通二进制状态。

\subsection{状态与组合计算为什么必须分开}

假设加法器的两个输入分别为 \(7\) 和 \(-2\)。只要输入稳定，加法器的输出就会稳定为
\(5\)。但加法器本身并不记得这个结果；输入变化后，输出随之变化。若下一条指令还要使用
\(5\)，就必须在规定时刻把它写入寄存器。寄存器在一个时钟边沿接收新值，在下一个边沿
到来前保持该值，这使连续动作能够首尾相接。

\begin{center}
\begin{tikzpicture}[node distance=12mm and 19mm]
\node[stage,minimum width=3.2cm] (old) {周期开始\\\code{r4=7}};
\node[stage,minimum width=3.7cm,right=of old] (alu) {组合加法\\\code{7+(-2)=5}};
\node[stage,minimum width=3.2cm,right=of alu] (new) {时钟边沿后\\\code{r4=5}};
\draw[flow] (old) -- node[above,font=\footnotesize]{读旧状态} (alu);
\draw[flow] (alu) -- node[above,font=\footnotesize]{提交新状态} (new);
\end{tikzpicture}
\end{center}
\diagramnote{图中的加法结果在时钟边沿之前只是候选值；写入寄存器后，它才成为下一条指令
可以读取的处理器状态。}

“提交”在本章中表示一次处理器状态更新已经越过时钟边界。它不表示数据已经断电保存，也
不表示整个任务完成。后续讨论存储和设备时，还会出现不同层次的完成边界。

\subsection{给教学处理器规定一个足够小的软件视图}

真实处理器的寄存器和指令很多，本章只保留运行该任务所需的部分。教学处理器有八个 32 位
通用寄存器 \code{r0--r7}，另有 PC、SP 和零标志 \code{Z}。寄存器的用途由启动约定和
任务代码共同决定：

\begin{longtable}{|L{0.17\linewidth}|L{0.28\linewidth}|L{0.43\linewidth}|}
\hline
状态 & 本次任务中的初始含义 & 执行期间怎样变化 \\ \hline
\code{PC} & 任务入口 \code{0x00100000} & 顺序前进，分支和返回会改写它 \\ \hline
\code{SP} & 指向预置返回地址 & 调用时下降，返回时上升 \\ \hline
\code{r0} & 输入数组首地址 \code{0x00102000} & 每轮增加 4，指向下一个整数 \\ \hline
\code{r1} & 元素个数 4 & 每轮减 1，为 0 时循环结束 \\ \hline
\code{r2} & 结果地址 \code{0x00102020} & 本任务中保持不变 \\ \hline
\code{r4} & 部分和，初值 0 & 每轮加上一个数组元素 \\ \hline
\code{r5} & 当前读出的整数 & 每轮被下一次装入覆盖 \\ \hline
\code{Z} & 最近一次规定运算的结果是否为 0 & 供条件分支读取 \\ \hline
\end{longtable}

这张表同时说明了一个重要边界：处理器不会识别“数组首地址”或“元素个数”。当固件把
\code{0x00102000} 写进 \code{r0} 时，硬件只得到一个 32 位位串。只有任务按照约定把它
用作装入指令的基址，它才表现为数组地址。

\subsection{指令也必须以字节形式存在}

处理器不能执行写在纸上的“求和”二字。它读取固定格式的机器指令位串，根据其中的操作码
和操作数字段打开不同的数据通路。本章采用固定 4 字节的教学指令。第一个字节是操作码，
第二、第三个字节通常给出寄存器编号，最后一个字节可表示小范围立即数或地址偏移。分支
指令把后两个字节解释为相对位移。

\begin{longtable}{|L{0.18\linewidth}|L{0.30\linewidth}|L{0.40\linewidth}|}
\hline
指令形式 & 状态变化 & 本章中的用途 \\ \hline
\code{MOVI rd,imm} & \code{rd <- imm} & 把部分和设为 0 \\ \hline
\code{LOAD rd,[rb+d]} & \code{rd <- MEM32[rb+d]} & 从数组读取一个整数 \\ \hline
\code{STORE rs,[rb+d]} & \code{MEM32[rb+d] <- rs} & 写回最终结果 \\ \hline
\code{ADD rd,rs} & \code{rd <- rd+rs} & 更新部分和 \\ \hline
\code{ADDI rd,imm} & \code{rd <- rd+imm} & 移动指针或减少计数 \\ \hline
\code{BNEZ rs,target} & \code{rs!=0} 时改写 PC & 控制循环 \\ \hline
\code{CALL target} & 压入返回地址并改写 PC & 调用子程序 \\ \hline
\code{RET} & 从栈中取回 PC & 返回调用者 \\ \hline
\end{longtable}

这不是某个商业处理器的真实编码，而是一个字段完整、状态变化明确的教学接口。后文讨论的
操作系统原理不依赖某个特定指令集；使用固定编码的意义在于，每次“取指”和“执行”都能
对应到确切地址、确切字节和确切状态变化。

求和任务的核心指令如下。左侧地址按每条指令 4 字节递增：

\begin{lstlisting}
0x00100000  MOVI  r4, 0
0x00100004  LOAD  r5, [r0+0]
0x00100008  ADD   r4, r5
0x0010000c  ADDI  r0, 4
0x00100010  ADDI  r1, -1
0x00100014  BNEZ  r1, 0x00100004
0x00100018  STORE r4, [r2+0]
0x0010001c  RET
\end{lstlisting}

第一条指令把 \code{r4} 置零。循环体每次读取一个 32 位整数，更新部分和，将输入指针移到
下一个整数，再减少剩余数量。计数不为零时，PC 回到 \code{0x00100004}；第四轮后分支不
成立，任务写出结果并执行 \code{RET}。

\begin{problemframe}
\textbf{复算点。} 若 \code{r0=0x00102000}、\code{r1=4}，循环执行两轮后应有
\code{r0=0x00102008}、\code{r1=2}、\code{r4=5}。这里只计算程序语义；这些字节怎样进入
内存、处理器怎样读到它们，仍未解决。
\end{problemframe}

\section{为什么既需要断电保存的位置，也需要运行期可写的位置}

现在已有一段机器指令，但还没有说明它在上电前位于哪里。若所有字节都只放在 RAM 中，
断电后内容消失；下一次通电时，PC 即使得到一个固定初值，也读不到可靠指令。机器因此需要
一片非易失启动存储，保存最早执行的固件。

仅有启动存储仍不够。求和任务要更新栈、部分数据和装入后的任务区域，而本章把启动存储
视为运行期只读。机器还需要 RAM，为会变化的字节提供位置。两类存储都能响应地址读取，
但它们的保留时间和写入规则不同，不能用一句“内存”混在一起。

\subsection{启动存储的三层含义}

\textbf{硬件模型。} 启动存储是一组按地址索引的非易失字节单元。输入是读地址和读请求，
输出是相应字节及完成状态。上电期间的普通写请求不会改变它；如何烧写固件属于机器制造或
维护阶段，本章不展开。

\textbf{软件模型。} 启动存储占用物理地址 \code{0x00000000} 至
\code{0x0000ffff}。复位入口、装入程序和返回入口都位于其中。程序使用普通取指或读指令访问这段
地址，但写入会得到错误响应。

\textbf{抽象。} 启动软件可以依赖一段跨断电保留且复位后立即可读的字节序列。这个抽象
只保证地址与字节，不提供文件名、目录或版本选择；这些含义若有需要，必须由固件格式另行
建立。

\subsection{RAM 的三层含义}

\textbf{硬件模型。} RAM 是按地址索引的易失存储单元阵列。读请求返回旧值，写请求按照
地址、写数据和字节使能更新选中单元。电源失效后，内容不再可靠。

\textbf{软件模型。} 物理地址 \code{0x00100000--0x001fffff} 可读写。任务代码、输入
数据、栈和固件工作区都必须在这一范围内划出不重叠区域。RAM 不保存“代码”“整数”或
“栈帧”标签，同一组字节的解释完全由访问它的指令和软件约定决定。

\textbf{抽象。} 软件得到一片可变的、按字节寻址的工作空间。当前抽象不包含虚拟地址、
访问权限和所有者；任何能够发出物理写请求的代码都可能改写其中任意位置。

\subsection{四个整数在内存中究竟是什么}

输入数组从 \code{0x00102000} 开始，每个元素占 4 字节。按低有效字节在前的次序，整数
\(7\) 存成 \code{07 00 00 00}，整数 \(-2\) 的 32 位补码为
\code{fe ff ff ff}。四个元素的实际布局如下：

\begin{longtable}{|L{0.22\linewidth}|L{0.25\linewidth}|L{0.41\linewidth}|}
\hline
地址范围 & 字节内容 & 按 32 位有符号整数解释 \\ \hline
\code{0x00102000--03} & \code{07 00 00 00} & \(7\) \\ \hline
\code{0x00102004--07} & \code{fe ff ff ff} & \(-2\) \\ \hline
\code{0x00102008--0b} & \code{0b 00 00 00} & \(11\) \\ \hline
\code{0x0010200c--0f} & \code{05 00 00 00} & \(5\) \\ \hline
\code{0x00102020--23} & 装入时清零 & 任务完成后应为 \code{15 00 00 00} \\ \hline
\end{longtable}

\code{LOAD r5,[r0+0]} 并不是“读取数组元素”这一高级动作。处理器先计算
\code{r0+0} 得到物理地址，再请求从该地址开始的 4 个字节，最后按规定字节序把它们组合
成 32 位值写入 \code{r5}。数组是连续元素的约定，装入指令只看到地址和宽度。

\subsection{存储容量与地址范围不是同一个概念}

本样机的 RAM 容量为 \(1\,\mathrm{MiB}\)，需要 \(2^{20}\) 个字节位置。处理器却能表达
\(2^{32}\) 个不同的物理地址。没有必要让每个地址都对应 RAM；地址还要留给启动存储和
设备。地址是请求中的编号，容量是某个器件实际实现的单元数量。二者通过后面引入的地址
互连建立对应关系。

到这里，机器已经有了保存启动字节和运行字节的器件，但处理器还没有连接到它们。下一步
必须说明一次读写请求携带什么，以及多个器件如何保证只有一个作出响应。
